\section{Exact arithmetic of the pair policies and of the ladder}\label{S:pairs}

All quantities are for the instance of Section~5 of the main text: $\ddot p=u$, $u\in U$ (the hexagon with vertices $\pm n_1,\pm n_2,\pm n_3$, $h_U(\pm n_j)=1$), $p_j(0)=\tfr{34}{25}n_j$, $v_j(0)=n_j$, facets $n_i^\top p\le2$ and $|v_i|\le\tfr65$, $\tau=\tfr15$, $T=\tfr65$. Along branch $j$ write $\zeta_j(t)=n_j^\top p_j(t)$ and $\nu_j(t)=n_j^\top v_j(t)$ for the critical coordinate and its velocity. Under a constant control $w$ on an interval of length $\ell$ starting from $(\zeta,\nu)$, $\zeta\mapsto\zeta+\nu\ell+\tfrac12(n_j^\top w)\ell^2$ and $\nu\mapsto\nu+(n_j^\top w)\ell$; braking at $u=-n_j$ from velocity $\nu>0$ to rest adds $\nu^2/2$ to $\zeta$ and takes time $\nu$.

\subsection{Proposition 5.2: the two pair policies at $\tau=\tfr15$}

\paragraph{Pair $\{1,2\}$, blind control $w=(-\tfr25,-\tfr45)$.}
$n_1^\top w=n_2^\top w=-\tfr25$ (both in-belief branches brake at the shared authority $\tfr25$). Window exit, both branches: $\nu=1-\tfr25\cdot\tfr15=\tfr{23}{25}$, $\zeta=\tfr{34}{25}+\tfr15-\tfrac12\cdot\tfr25\cdot\tfr1{25}=\tfr{194}{125}$. Braking to rest: peak $\zeta=\tfr{194}{125}+\tfrac12(\tfr{23}{25})^2=\tfr{1940+529}{1250}=\tfr{2469}{1250}=1.9752$, rest at $t=\tfr15+\tfr{23}{25}=\tfr{28}{25}<T$; slack $\tfr{31}{1250}$. Velocity components at the junctions: branch~1: $(1,0)\to(\tfr{23}{25},-\tfr4{25})\to(0,-\tfr4{25})$; branch~2: $(-\tfr35,\tfr45)\to(-\tfr{17}{25},\tfr{16}{25})\to(-\tfr{16}{125},-\tfr{12}{125})$. Every component stays in $[-1,1]$, so the velocity facets $|v_i|\le\tfr65$ hold with margin at least $\tfr15$. Noncritical position facets: for $i\ne j$, $n_i^\top p_j(t)\le(\tfr{34}{25}+t)\,n_i^\top n_j+\tfrac12t^2\le\tfr{18}{25}$ on $[0,T]$ since $n_i^\top n_j\le-\tfr7{25}$ and $h_U(n_i)=1$. The pair $\{1,3\}$ is the reflection $v_2\mapsto-v_2$.

\paragraph{Pair $\{2,3\}$, blind control $w=n_1=(1,0)$.}
$n_2^\top n_1=n_3^\top n_1=-\tfr35$. Window exit: $\nu=1-\tfr35\cdot\tfr15=\tfr{22}{25}$, $\zeta=\tfr{34}{25}+\tfr15-\tfrac12\cdot\tfr35\cdot\tfr1{25}=\tfr{387}{250}$. Braking: peak $\tfr{387}{250}+\tfrac12(\tfr{22}{25})^2=\tfr{1935+484}{1250}=\tfr{2419}{1250}=1.9352$, rest at $t=\tfr{27}{25}$; slack $\tfr{81}{1250}$. Velocity components, branch~2: $(-\tfr35,\tfr45)\to(-\tfr25,\tfr45)\to(-\tfr25,\tfr45)+\tfr{22}{25}(\tfr35,-\tfr45)=(\tfr{16}{125},\tfr{12}{125})$; branch~3 is the mirror image. All components have magnitude at most $\tfr45$.

\paragraph{Falsified variants (recorded so that the policies are not read as generic).}
(i) Holding $u=0$ through the blind window and then braking at $-n_j$ peaks at $\tfr{34}{25}+\tfr15+\tfr12=\tfr{103}{50}>2$ on every branch. (ii) For $\{2,3\}$, inserting a reduced-brake leg $u=-\tfr35n_j$ of length $\tfr15$ after the window before full braking gives $(\zeta,\nu)=(\tfr{387}{250}+\tfr{22}{25}\cdot\tfr15-\tfrac12\cdot\tfr35\cdot\tfr1{25},\ \tfr{22}{25}-\tfr3{25})=(\tfr{214}{125},\tfr{19}{25})$ at $t=\tfr25$ and then peaks at $\tfr{214}{125}+\tfrac12(\tfr{19}{25})^2=\tfr{2501}{1250}>2$ at $t=\tfr{29}{25}$. (iii) The \emph{off-belief} branch~1 under the $\{2,3\}$ blind control $u=n_1$ accelerates to $v_1(\tau)=\tfr65$, exactly the velocity bound, and would then peak at $\tfr{34}{25}+\tfr15+\tfr1{50}+\tfr12(\tfr65)^2=\tfr{23}{10}$: the pair policy is a policy for the pair, not for the triple.

\subsection{Theorem 6.2: junction values along the threshold schedules}

Let $t=t^\ast(B)$, $q=1-t\tau$, and the schedule ``$w$ on $[0,\tau]$; $-n_j$ on $[\tau,\tau+q]$; $0$ on $[\tau+q,T]$''. The critical coordinate obeys
\[
\zeta_j(\tau)=\tfr{34}{25}+\tau-\tfrac t2\tau^2,\qquad\nu_j(\tau)=1-t\tau=q,\qquad\zeta^{\max}=\zeta_j(\tau)+\tfrac12q^2=\tfr{93}{50}+(1-t)\tau+\tfrac{t^2-t}2\tau^2 ,
\]
and $\zeta^{\max}=2$ is the quadratic $6\tau^2-30\tau+7=0$ for $t=\tfr25$ and $6\tau^2-20\tau+7=0$ for $t=\tfr35$ (multiply $2-\zeta^{\max}$ by $50$). The smaller roots are $\tau_{12}=\tfr{15-\sqrt{183}}6$ and $\tau_{23}=\tfr{10-\sqrt{58}}6$, with rational brackets $0.2453<\tau_{12}<0.2454$ and $0.3973<\tau_{23}<0.3974$ obtained by sign changes of the exact quadratics (the verifier evaluates them at the bracket endpoints in rational arithmetic). The label time $s_\ast=\tau+q=1+(1-t)\tau$ is $<\tfr65$ at both roots ($1+\tfr35\cdot0.2454<1.148$; $1+\tfr25\cdot0.3974<1.159$).

\emph{Velocity junction values at the roots.} Each velocity coordinate is affine on each leg, so its extreme values occur at $t=0$, $\tau$, $\tau+q$ and $T$. For $B=\{1,2\}$ with $w=(-\tfr25,-\tfr45)$ at $\tau=\tau_{12}$: branch~1, $(1,0)\to(1-\tfr25\tau,-\tfr45\tau)\to\tau(w+\tfr25n_1)=(0,-\tfr45\tau)$; branch~2, $(-\tfr35,\tfr45)\to(-\tfr35-\tfr25\tau,\tfr45-\tfr45\tau)\to\tau(w+\tfr25n_2)=\tau(-\tfr{16}{25},-\tfr{12}{25})$. For $B=\{2,3\}$ with $w=(1,0)$ at $\tau=\tau_{23}$: branch~2, $(-\tfr35,\tfr45)\to(-\tfr35+\tau,\tfr45)\to\tau(w+\tfr35n_2)=\tau(\tfr{16}{25},\tfr{12}{25})$, and branch~3 its mirror image. With $\tau<\tfr25$ every entry has magnitude at most $1$, so $|v_i|\le1<\tfr65$ throughout, including the final hold leg where the velocity is constant. The noncritical position bound $\tfr{18}{25}$ of the previous subsection is delay-independent.

\emph{Monotonicity of the schedule in $\tau$.} $\zeta^{\max}(\tau)$ is a quadratic in $\tau$ with vertex at $\tau=1/t$, hence increasing on $[0,1/t)\supset[0,\tfr12]$; so the schedule is safe for every $\tau\le\tau_B$, and at $\tau=\tau_B$ exactly it touches the ceiling without crossing it (constraints are closed).

\subsection{Proposition 6.1(ii): the shared authorities}

For a pair $B=\{i,j\}$ the objective $\min_{w\in U}\max(n_i^\top w,n_j^\top w)$ is piecewise linear; its minimum over the hexagon is attained at a vertex of $U$ or on the bisector $\{(n_i-n_j)^\top w=0\}$. For $\{1,2\}$: $n_1-n_2=(\tfr85,-\tfr45)$, bisector $w_2=2w_1$; on it $n_1^\top w=w_1$, and the hexagon's extreme point with $w_1$ minimal is where $|w_2|\le\tfr45$ binds, $w=(-\tfr25,-\tfr45)$, giving $-\tfr25$. At the six vertices $\max(n_1^\top w,n_2^\top w)$ takes the values $1,\tfr35,1,\tfr35,-\tfr7{25},\tfr35$ (at $n_1,-n_1,n_2,-n_2,n_3,-n_3$), none below $-\tfr7{25}$; hence $t^\ast(\{1,2\})=\tfr25$. For $\{2,3\}$: bisector $w_2=0$, $n_2^\top w=-\tfr35w_1$, minimal at the vertex $w=(1,0)=n_1$: $t^\ast=\tfr35$. For the triple, $\sum_j\lambda_jn_j=0$ forces $\max_jn_j^\top w\ge0$ for every $w$, with equality at $w=0$. The pooled normals of the pair duals, $\tfr38n_1+\tfr58n_2=(0,\tfr12)$, $\tfr38n_1+\tfr58n_3=(0,-\tfr12)$ and $\tfr12n_2+\tfr12n_3=(-\tfr35,0)$, have support values $h_U(-\nu)$ equal to $\tfr25$, $\tfr25$, $\tfr35$ respectively ($h_U(e_2)=\tfr45$, $h_U(e_1)=1$)---the same numbers, which is why the two-label duals are tight against the schedules.
